\section{Methodology}\label{sec:methodology}

\subsection{Solar geometry and azimuth convention}\label{subsec:geometry}
Let $\omega$ denote the solar hour angle and $\beta$ the surface tilt. We
define $\gamma=0$ for an equator-facing surface, with the sign convention
made explicit for the hemisphere under study. The daily hour angle of maximum
clear-sky GTI is denoted by $\omega^{*}$.
\draftnote{Add solar-angle equations, the exact implementation of solar time and timestamps, timezone convention, and the explicit signed azimuth definition (east/west) used by Azimut2026. Note how the convention maps to pvlib if applicable.}

\begin{figure}[t]
\centering
\figplaceholder{Geometry: tilt $\beta$, azimuth $\gamma$, hour angle $\omega$, and the controlled mount.}
\caption{Conceptual geometry and experimental configuration. Replace the
placeholder with a diagram and/or a photograph of the measurement system.}
\label{fig:geometry}
\end{figure}

\subsection{Mapping the hour angle of maximum GTI to surface azimuth}
\label{subsec:mapping}
For a selected day $d$, the modeled relationship between surface azimuth and
the location of the GTI maximum is represented by
\begin{equation}
  \hat\gamma_d = \mathcal{F}_d\!\left(\hat\omega_d^{*};\beta,\phi,\boldsymbol{p}_d\right),
  \label{eq:azimuth_mapping}
\end{equation}
where $\phi$ is the site latitude and $\boldsymbol{p}_d$ denotes the
assumptions and inputs used to construct the clear-sky/transposition mapping.
The 2021 formulation used a fitted mapping between $\omega^{*}$ and $\gamma$
\citep{piccioli2021}.
\draftnote{Describe the mapping ACTUALLY implemented in the repository (analytical fit, LUT, interpolator, clear-sky inputs, transposition model). Do not assume that the 2021 tangent fit is still used. Specify the azimuth search domain and numerical resolution.}

\subsection{Estimation of the daily GTI maximum}\label{subsec:dailyfit}
For a suitable interval around the daily maximum, a quadratic approximation
can be written as
\begin{equation}
  \GTI(\omega)\simeq a_2\omega^2+a_1\omega+a_0,
  \qquad
  \hat\omega^{*}=-\frac{a_1}{2a_2},\qquad a_2<0.
  \label{eq:quadratic}
\end{equation}
The location of the fitted maximum, rather than the largest individual
measurement, is used to reduce sensitivity to high-frequency fluctuations.
\draftnote{Specify fit-window definition, sample cadence, fitting algorithm, iteration/stopping rules, physical constraints, and treatment of non-concave or poorly constrained fits. Verify these against current code.}

\subsection{Clear-sky selection, outliers, and diagnostic indicators}
\label{subsec:quality}
A daily estimate is only meaningful if the irradiance maximum is sampled
adequately and the retained GTI observations are suitable for fitting.
Candidate diagnostics include fit residual dispersion, fraction of rejected
points, residual asymmetry, maximum curvature, and temporal data coverage.
\draftnote{Define the actual clear-sky and outlier masks with precise thresholds. Test the exceptional dates 2026-07-01, 2026-08-31, and 2026-09-17; distinguish physical cloud variability from an algorithmic selection artifact. Do not tune screening criteria solely to minimize observed angular error.}

\subsection{Uncertainty propagation and confidence assessment}
\label{subsec:uncertainty}
Writing $\mathbf{a}=(a_2,a_1,a_0)$, the formal local uncertainty of
$\hat\omega^{*}$ may be approximated by the delta method,
\begin{equation}
  u^2_{\omega^{*}}\simeq
  \left(\nabla_{\mathbf{a}}\omega^{*}\right)^{\mathsf{T}}
  \boldsymbol{\Sigma}_{\mathbf{a}}
  \nabla_{\mathbf{a}}\omega^{*},
  \label{eq:omega_uncertainty}
\end{equation}
with $\boldsymbol{\Sigma}_{\mathbf{a}}$ the covariance matrix of the
fitted coefficients. A first-order propagated component is
\begin{equation}
 u_{\gamma,\mathrm{fit}} \simeq
 \left|\frac{\partial\mathcal{F}_d}{\partial\omega^{*}}\right|
 u_{\omega^{*}}.
 \label{eq:gamma_uncertainty}
\end{equation}
This formal uncertainty does not, by itself, account for correlated residuals,
incorrect atmospheric inputs, timestamp offsets, or errors in the reference
orientation. Alternative uncertainty estimates and their calibration should
therefore be evaluated.
\draftnote{Compare covariance propagation with a time-aware bootstrap (e.g., blocks) and with empirical between-day spread. Report uncertainty components separately.}

\subsection{Multi-day estimation and temporal stability}
\label{subsec:aggregation}
For a period expected to have stable orientation, daily estimates
$\hat\gamma_d$ can be combined using robust or uncertainty-weighted procedures.
A weighted mean may be considered as a baseline,
\begin{equation}
  \hat\gamma_{\mathrm{period}}=
  \frac{\sum_d w_d\hat\gamma_d}{\sum_d w_d},
  \qquad w_d=u_{\gamma,d}^{-2},
  \label{eq:aggregation}
\end{equation}
provided that the daily uncertainties are calibrated and the effects of
correlations and shared systematic errors are addressed.
\draftnote{Compare median, robust location estimator, and weighted mean. Specify change-point detection procedure for historical stations, if used.}

\subsection{Validation protocol and evaluation metrics}
\label{subsec:metrics}
The azimuth error for controlled orientation $j$ is
$\Delta\gamma_j=\hat\gamma_j-\gamma_{\mathrm{ref},j}$.
Performance is described by bias, mean absolute error, and root mean square
error, together with repeatability across days:
\begin{equation}
  \MBE_{\gamma}=\frac{1}{N}\sum_{j=1}^{N}\Delta\gamma_j,\quad
  \MAE_{\gamma}=\frac{1}{N}\sum_{j=1}^{N}|\Delta\gamma_j|,\quad
  \RMSE_{\gamma}=\sqrt{\frac{1}{N}\sum_{j=1}^{N}\Delta\gamma_j^{2}}.
  \label{eq:metrics}
\end{equation}
\draftnote{Define unit of analysis (orientation period vs day), between-day dependence, confidence intervals, and whether reference orientation uncertainty is included. Establish an untouched test set or cross-period protocol before screening/model selection.}
